\documentclass{scrlttr2}
\usepackage[utf8]{inputenc}
\usepackage[T1]{fontenc}
\usepackage{geometry}
\usepackage{graphicx}
\usepackage{float}
\usepackage{hyperref}
\usepackage[english,indonesian]{babel}

\title{Optimasi Pemrosesan Template LaTeX Menggunakan Pendekatan Berbasis AI}
\author{Academic Support Agent}
\date{\today}

\begin{document}
\maketitle

\begin{abstract}

\end{abstract}

\chapter{PENDAHULUAN}


\chapter{TINJAUAN PUSTAKA}


% Original content (commented out):
% \documentclass[]{scrlttr2}
\usepackage[T1]{fontenc}
\usepackage[utf8]{inputenc}
\usepackage[ngerman]{babel}
\usepackage{geometry}
\geometry{verbose,a4paper,tmargin=25mm,bmargin=25mm,lmargin=20mm,rmargin=20mm}
\usepackage{graphicx}

\begin{document}
	
\pagestyle{empty}
% If you want headings on subsequent pages,
% remove the ``%'' on the next line:
% \pagestyle{headings}
\setkomavar{fromname}{Marie Mustermann}
\setkomavar{fromaddress}{Musterstraße 2\\0002 Musterstadt-Neustadt}
\setkomavar{signature}{(Marie Mustermann)}

\setkomavar{subject}{}

\begin{letter}{ 
	% address of recipient
    Max Mustermann 	\\
    Musterstraße 1	\\
    0001 Musterstadt\\
}

\opening{Sehr geehrte Dame, sehr geehrter Herr,}

Dies ist ein Blindtext.

\closing{Mit freundlichen Grüßen}

%\includegraphics[height=30px]{ %Unterschrift}

\ps{}

%enclosure listing
%\encl{}

\end{letter}
\end{document}


\end{document}